\chapter{Introduction}

\section{Motivation}

State-of-the-art (SOTA) object classification methods utilize deep learning models that are either trained end-to-end on the target application or initially trained in one application and then fine-tuned for a different target application. In either case, they are monolithic models: the same single model is normally trained on a vast dataset and used for feature extraction and classification. These models are highly effective when the set of object classes to be recognized is fixed and known at the training/fine-tuning time. However, when new object classes need to be incorporated, the model must be retrained or fine-tuned. This training requires significant computational resources and access to the entire training dataset and, in some cases, leads to catastrophic forgetting \cite{catastrophic}.

One way to integrate new classes while avoiding catastrophic forgetting, is to approach the problem using continual or class-incremental object classification, where an algorithm must learn to distinguish between a growing number of object classes \cite{vandeVen}. In the context of continual learning, a task refers to a distinct portion of the training experience, usually presented sequentially to the model. Instead of having access to the entire dataset at once, as in standard supervised learning, the model receives data in sequential chunks \cite{taskarithmeticloracontinual}. Each task generally correspond to a set of new classes to recognize. The main challenge is that when the model is trained task by task, it tends to overwrite previously learned knowledge. Given that current monolithic models are not capable of handling this type of application, an alternative paradigm is required.

This limitation has motivated SOTA methods for object classification to focus on balancing stability and plasticity. Stability refers to a model’s ability to preserve previously learned knowledge, while plasticity describes its capacity to adapt and incorporate new tasks \cite{API}. However, traditional approaches rarely optimize both aspects simultaneously: when stability dominates, the model struggles to integrate new knowledge; when plasticity dominates, it loses access to past information. This trade-off is compounded by the limited capacity of neural networks to integrate new representations without interfering with old ones, and by the need to continually adjust the network’s structure to match the increasing number of object classes.

 One solution to this problem is to divide the monolithic models into two distinct components that must be trained separately, at different times: one responsible for feature extraction and the other for object classification. The feature extraction model must be trained when the recognition system is initialized/created and the obtained features must be general enough to characterize all possible classes of objects to be seen (not just the ones known at training time). The object classification model should be easily updated or trained as new object classes are detected. Ideally, it should be able to incorporate new object classes without modifying the parts of the model responsible for recognizing the already known object classes. This mirrors the functionality of the human visual system, where the structures responsible for the extraction of characteristics (simple and complex cells in cortical areas V1 and V2) undergo adaptation during the initial postnatal weeks/months \cite{Sluyters90}. However, as object recognition occurs in high-level cognitive areas, and acquiring knowledge of a new object class can result in a slight reduction in the ability to recognize similar classes of objects over time, the ability to learn new categories of objects persists throughout life. 

We are inspired by this idea and propose a continual learning object recognition system based on two components. The first component is a feature extraction model built using a pre-trained foundation model. The second component is a statistical classifier built using various Gaussian models. Each model characterizes a specific object class independently of the others. Since these Gaussian models can be trained or built at different times, a new Gaussian model is created to characterize a new object class when it needs to be incorporated. For classification, the class whose Gaussian model yields the highest likelihood of the input features is selected. This approach is analogous to a Bayes classifier, in which new classes are incorporated over time.

Gaussian models offer two key advantages for characterizing object classes: first, each object class model operates independently from the others, facilitating the straightforward incorporation of new classes into the model. Second, Gaussian models provide a compact representation of object classes, requiring few samples to construct a model that can be refined as new samples become available.

\section{Hypothesis}

We hypothesize that the combination of independent Gaussian models and a pretrained feature extractor provides a robust framework for continual object classification. By enabling the addition of new classes without full-model retraining, this design eliminates catastrophic forgetting. The inherent simplicity of this approach is a key advantage, yielding performance that rivals state-of-the-art methods while maintaining lower computational overhead.



\section{General Objectives}
The main objective of this work is to develop and validate a continual learning framework for object classification, using pre-trained foundation models as features extractors and independent Gaussian models, with the purpose of enabling the incremental incorporation of new classes while preventing catastrophic forgetting of previously learned class.


\section{Specific Objectives}

For the purpose of fulfilling the general objective and verifying the hypothesis of this work, the following specific objectives are defined:

\begin{itemize}
	
	\item To assess the Gaussianity of feature distributions within the latent space to validate the underlying distributional assumptions.

	
	\item To characterize the efficacy of independent Gaussian distributions as class representations within a continual learning framework.
	
	\item To quantify the impact of various covariance matrix approximation methods on model stability and classification accuracy.
	
	\item To evaluate the generalization capabilities of diverse feature extractors across multiple dataset configurations and incremental learning scenario.

			
\end{itemize}

\section{Thesis Structure}


